\chapter{Resultados}\label{cap:resultados}

Todas las cifras de este capítulo salen del tercer run del kernel y se evalúan en los sujetos de prueba de cada fold.
Los intervalos son del 95\pct{}, por bootstrap de sujetos.

\section{Un sujeto para empezar}

Antes de los promedios, un caso. La \cref{fig:ilustracion} muestra el peso de un sujeto real de la campaña~3 que en el
fold~0 es de prueba. Antes del reposo, su peso oscila alrededor de la media del pre-reposo. En cama baja, como en casi
todos, y se queda abajo.

\begin{figure}[htbp]
\begin{tikzpicture}
\begin{axis}[informe, width=\linewidth, height=6.6cm, xmin=-15, xmax=90, ymin=84.14, ymax=91.24, ymajorgrids, xlabel={día respecto del inicio del reposo en cama}, ylabel={peso (kg)}, axis x line*=bottom, axis y line*=left, clip=false, legend style={at={(0.99,0.99)}, anchor=north east, fill=white, fill opacity=0.85, text opacity=1}, legend columns=1]
\fill[fondo] (axis cs:-15,84.14) rectangle (axis cs:-0.5,91.24);
\node[anchor=north west, font=\scriptsize\sffamily\color{tinta2}] at (axis cs:-14.5,91.24) {pre-reposo};
\fill[serieB, opacity=0.13] (axis cs:-0.5,87.367) rectangle (axis cs:90,89.826);
\draw[serieB, line width=1pt, densely dashed] (axis cs:-0.5,88.596) -- (axis cs:90,88.596);
\fill[serieA, opacity=0.18] plot coordinates {(0,90.844) (1,90.561) (2,90.393) (3,90.273) (4,90.179) (5,90.102) (6,90.036) (7,89.978) (8,89.926) (9,89.879) (10,89.837) (11,89.798) (12,89.761) (13,89.727) (14,89.695) (15,89.665) (16,89.636) (17,89.609) (18,89.583) (19,89.558) (20,89.535) (21,89.512) (22,89.489) (23,89.468) (24,89.448) (25,89.428) (26,89.408) (27,89.389) (28,89.371) (29,89.353) (30,89.336) (31,89.319) (32,89.302) (33,89.286) (34,89.270) (35,89.254) (36,89.239) (37,89.224) (38,89.210) (39,89.195) (40,89.181) (41,89.167) (42,89.153) (43,89.140) (44,89.127) (45,89.114) (46,89.101) (47,89.088) (48,89.076) (49,89.063) (50,89.051) (51,89.039) (52,89.027) (53,89.015) (54,89.004) (55,88.992) (56,88.981) (57,88.970) (58,88.959) (59,88.948) (60,88.937) (61,88.926) (62,88.915) (63,88.905) (64,88.894) (65,88.884) (66,88.874) (67,88.864) (68,88.853) (69,88.843) (70,88.834) (71,88.824) (72,88.814) (73,88.804) (74,88.795) (75,88.785) (76,88.776) (77,88.766) (78,88.757) (79,88.748) (80,88.738) (81,88.729) (82,88.720) (83,88.711) (84,88.702) (85,88.693) (86,88.684) (87,88.675) (88,88.667) (89,88.658) (89,84.540) (88,84.549) (87,84.558) (86,84.566) (85,84.575) (84,84.584) (83,84.593) (82,84.602) (81,84.611) (80,84.620) (79,84.630) (78,84.639) (77,84.648) (76,84.658) (75,84.667) (74,84.677) (73,84.686) (72,84.696) (71,84.706) (70,84.716) (69,84.726) (68,84.736) (67,84.746) (66,84.756) (65,84.766) (64,84.776) (63,84.787) (62,84.797) (61,84.808) (60,84.819) (59,84.830) (58,84.841) (57,84.852) (56,84.863) (55,84.874) (54,84.886) (53,84.897) (52,84.909) (51,84.921) (50,84.933) (49,84.945) (48,84.958) (47,84.970) (46,84.983) (45,84.996) (44,85.009) (43,85.022) (42,85.036) (41,85.049) (40,85.063) (39,85.077) (38,85.092) (37,85.106) (36,85.121) (35,85.137) (34,85.152) (33,85.168) (32,85.184) (31,85.201) (30,85.218) (29,85.235) (28,85.253) (27,85.271) (26,85.290) (25,85.310) (24,85.330) (23,85.350) (22,85.372) (21,85.394) (20,85.417) (19,85.440) (18,85.465) (17,85.491) (16,85.519) (15,85.547) (14,85.577) (13,85.609) (12,85.643) (11,85.680) (10,85.719) (9,85.762) (8,85.808) (7,85.860) (6,85.918) (5,85.984) (4,86.061) (3,86.155) (2,86.275) (1,86.443) (0,86.726)} -- cycle;
\draw[tinta3, densely dashed] (axis cs:28,84.14) -- (axis cs:28,91.24) node[pos=0.02, anchor=south west, font=\scriptsize\sffamily\color{tinta2}] {desde aquí se evalúa};
\addlegendimage{area legend, fill=serieB, fill opacity=0.25, draw=serieB}\addlegendentry{rango de Tierra: media del pre-reposo $\pm 1{,}96\,$DE}
\addlegendimage{area legend, fill=serieA, fill opacity=0.3, draw=serieA}\addlegendentry{rango adaptado ($k=28$, ancho calibrado)}
\addlegendimage{only marks, mark=*, mark size=1.6pt, tinta}\addlegendentry{peso medido cada mañana}
\addplot[serieA, line width=1.2pt, forget plot] coordinates {(0,88.785) (1,88.502) (2,88.334) (3,88.214) (4,88.120) (5,88.043) (6,87.977) (7,87.919) (8,87.867) (9,87.820) (10,87.778) (11,87.739) (12,87.702) (13,87.668) (14,87.636) (15,87.606) (16,87.577) (17,87.550) (18,87.524) (19,87.499) (20,87.476) (21,87.453) (22,87.431) (23,87.409) (24,87.389) (25,87.369) (26,87.349) (27,87.330) (28,87.312) (29,87.294) (30,87.277) (31,87.260) (32,87.243) (33,87.227) (34,87.211) (35,87.195) (36,87.180) (37,87.165) (38,87.151) (39,87.136) (40,87.122) (41,87.108) (42,87.094) (43,87.081) (44,87.068) (45,87.055) (46,87.042) (47,87.029) (48,87.017) (49,87.004) (50,86.992) (51,86.980) (52,86.968) (53,86.956) (54,86.945) (55,86.933) (56,86.922) (57,86.911) (58,86.900) (59,86.889) (60,86.878) (61,86.867) (62,86.856) (63,86.846) (64,86.835) (65,86.825) (66,86.815) (67,86.805) (68,86.795) (69,86.785) (70,86.775) (71,86.765) (72,86.755) (73,86.745) (74,86.736) (75,86.726) (76,86.717) (77,86.707) (78,86.698) (79,86.689) (80,86.679) (81,86.670) (82,86.661) (83,86.652) (84,86.643) (85,86.634) (86,86.625) (87,86.617) (88,86.608) (89,86.599)};
\addplot[only marks, mark=*, mark size=1.5pt, tinta, forget plot] coordinates {(-14,87.750) (-13,88.150) (-12,88.000) (-11,88.000) (-10,88.800) (-9,88.850) (-8,88.150) (-7,88.400) (-6,88.050) (-5,88.750) (-4,88.950) (-3,89.650) (-2,89.750) (-1,89.100) (0,89.900) (1,88.300) (2,88.800) (3,88.350) (4,88.270) (5,88.200) (6,88.150) (7,88.100) (8,88.100) (9,87.650) (10,87.750) (11,87.800) (12,87.750) (13,88.100) (14,87.650) (15,87.500) (16,87.200) (17,87.300) (18,87.300) (19,87.450) (20,86.600) (21,87.050) (22,87.350) (23,87.400) (24,87.350) (25,87.850) (26,86.600) (27,86.500) (28,86.900) (29,86.300) (30,86.650) (31,87.450) (32,87.050) (33,86.900) (34,87.500) (35,87.350) (36,87.050) (37,86.900) (38,86.800) (39,87.600) (40,86.700) (41,86.500) (42,86.800) (43,86.350) (44,87.050) (45,87.200) (46,87.750) (47,87.150) (48,86.700) (49,86.650) (50,87.200) (51,86.750) (52,86.850) (53,87.300) (54,87.000) (55,87.150) (56,87.250) (57,87.450) (58,87.800) (59,86.450) (60,86.950) (61,87.100) (62,86.850) (63,87.450) (64,87.000) (65,86.750) (66,87.150) (67,87.100) (68,87.350) (69,87.200) (70,86.950) (71,86.800) (72,87.200) (73,87.350) (74,87.900) (75,87.450) (76,87.150) (77,87.200) (78,87.350) (79,87.350) (80,87.350) (81,87.300) (82,87.250) (83,87.900) (84,87.200) (85,87.550) (86,87.650) (87,87.400) (88,87.300) (89,87.970)};
\end{axis}
\end{tikzpicture}

\caption[Un sujeto real: rango de Tierra y rango adaptado]{Peso de un sujeto real (C3\_6043, fold~0, sujeto de
prueba). El rango de Tierra (media del pre-reposo $\pm$\cifra{ilus.sd}\un{kg}) deja fuera
\cifra{ilus.fuera.tierra} de sus \cifra{ilus.dias} días evaluados desde el día 28. El rango adaptado, con sus
primeros 28 días en cama y un ancho calibrado de $\pm$\cifra{ilus.ancho}\un{kg}, no deja fuera ninguno
(\cifra{ilus.fuera.adapt}).}
\label{fig:ilustracion}
\end{figure}

El rango adaptado de este sujeto es más ancho que el de Tierra, porque el ancho calibrado es uno para toda la
población. Lo que cambia es el centro: sigue el nivel nuevo de la persona. En promedio, como se verá, el rango
adaptado también es más angosto que el de Tierra cuando a ambos se les exigen las mismas falsas alarmas.

\section{P1. El monitor de Tierra no sirve en cama}\label{sec:res-p1}

La \cref{fig:falsas} y la \cref{tab:rango} muestran cuántos días normales en cama quedan fuera del rango de Tierra. En el
peso son el \cifra{peso.tierra.fa}\pct{} \cifra{peso.tierra.fa.ic}: más de la mitad de los días, una persona sana recibiría una
alarma solo por estar en cama. En la presión diastólica, el pulso y la respiración son
\cifra{pa_dia.tierra.fa}, \cifra{pulso.tierra.fa} y \cifra{fr.tierra.fa}\pct. Promediando las siete variables sin
el agua, el monitor de Tierra da falsa alarma en el \cifra{prom.tierra.fa}\pct{} de los días.

\begin{figure}[htbp]
\begin{tikzpicture}
\begin{axis}[informe, width=0.86\linewidth, height=7cm, xmin=0, xmax=60, ymin=-0.6, ymax=7.6, ytick={0,...,7}, yticklabels={Agua,Temperatura,PA sistólica,Respiración,Diuresis,Pulso,PA diastólica,Peso},
 xmajorgrids, xtick={0,10,...,60}, xticklabel={\pgfmathprintnumber{\tick}\,\%},
 xlabel={días normales en cama fuera del rango (falsa alarma diaria)},
 y axis line style={draw=none}, ytick style={draw=none}, axis x line*=bottom, clip=false,
 legend style={at={(0.98,0.04)}, anchor=south east}, legend columns=1]
\addlegendimage{only marks, mark=*, mark size=2.8pt, serieB}\addlegendentry{monitor de Tierra}
\addlegendimage{only marks, mark=*, mark size=2.8pt, serieA}\addlegendentry{adaptado}
\addlegendimage{only marks, mark=|, mark size=4pt, tinta, line width=1.2pt}\addlegendentry{placebo dentro del pre-reposo}
\draw[tinta3, densely dashed] (axis cs:5,-0.6) -- (axis cs:5,7.6) node[pos=1, above, font=\scriptsize\sffamily\color{tinta2}] {5\,\%};
\draw[reja, line width=3pt] (axis cs:5.340,0) -- (axis cs:5.070,0);
\draw[tinta, line width=1.2pt] (axis cs:7.468,-0.28) -- (axis cs:7.468,0.28);
\fill[serieB, draw=white, line width=0.8pt] (axis cs:5.070,0) circle[radius=3pt];
\fill[serieA, draw=white, line width=0.8pt] (axis cs:5.340,0) circle[radius=3pt];
\node[anchor=west, font=\scriptsize\sffamily\color{tinta2}] at (axis cs:7.468,0) {\enspace5\,\%};
\draw[reja, line width=3pt] (axis cs:5.190,1) -- (axis cs:5.980,1);
\draw[tinta, line width=1.2pt] (axis cs:8.743,0.72) -- (axis cs:8.743,1.28);
\fill[serieB, draw=white, line width=0.8pt] (axis cs:5.980,1) circle[radius=3pt];
\fill[serieA, draw=white, line width=0.8pt] (axis cs:5.190,1) circle[radius=3pt];
\node[anchor=west, font=\scriptsize\sffamily\color{tinta2}] at (axis cs:8.743,1) {\enspace6\,\%};
\draw[reja, line width=3pt] (axis cs:4.820,2) -- (axis cs:10.960,2);
\draw[tinta, line width=1.2pt] (axis cs:8.081,1.72) -- (axis cs:8.081,2.2800000000000002);
\fill[serieB, draw=white, line width=0.8pt] (axis cs:10.960,2) circle[radius=3pt];
\fill[serieA, draw=white, line width=0.8pt] (axis cs:4.820,2) circle[radius=3pt];
\node[anchor=west, font=\scriptsize\sffamily\color{tinta2}] at (axis cs:10.960,2) {\enspace11\,\%};
\draw[reja, line width=3pt] (axis cs:4.980,3) -- (axis cs:13.770,3);
\draw[tinta, line width=1.2pt] (axis cs:12.152,2.7199999999999998) -- (axis cs:12.152,3.2800000000000002);
\fill[serieB, draw=white, line width=0.8pt] (axis cs:13.770,3) circle[radius=3pt];
\fill[serieA, draw=white, line width=0.8pt] (axis cs:4.980,3) circle[radius=3pt];
\node[anchor=west, font=\scriptsize\sffamily\color{tinta2}] at (axis cs:13.770,3) {\enspace14\,\%};
\draw[reja, line width=3pt] (axis cs:5.140,4) -- (axis cs:4.720,4);
\draw[tinta, line width=1.2pt] (axis cs:9.797,3.7199999999999998) -- (axis cs:9.797,4.28);
\fill[serieB, draw=white, line width=0.8pt] (axis cs:4.720,4) circle[radius=3pt];
\fill[serieA, draw=white, line width=0.8pt] (axis cs:5.140,4) circle[radius=3pt];
\node[anchor=west, font=\scriptsize\sffamily\color{tinta2}] at (axis cs:9.797,4) {\enspace5\,\%};
\draw[reja, line width=3pt] (axis cs:5.140,5) -- (axis cs:13.670,5);
\draw[tinta, line width=1.2pt] (axis cs:7.848,4.72) -- (axis cs:7.848,5.28);
\fill[serieB, draw=white, line width=0.8pt] (axis cs:13.670,5) circle[radius=3pt];
\fill[serieA, draw=white, line width=0.8pt] (axis cs:5.140,5) circle[radius=3pt];
\node[anchor=west, font=\scriptsize\sffamily\color{tinta2}] at (axis cs:13.670,5) {\enspace14\,\%};
\draw[reja, line width=3pt] (axis cs:5.410,6) -- (axis cs:14.910,6);
\draw[tinta, line width=1.2pt] (axis cs:9.091,5.72) -- (axis cs:9.091,6.28);
\fill[serieB, draw=white, line width=0.8pt] (axis cs:14.910,6) circle[radius=3pt];
\fill[serieA, draw=white, line width=0.8pt] (axis cs:5.410,6) circle[radius=3pt];
\node[anchor=west, font=\scriptsize\sffamily\color{tinta2}] at (axis cs:14.910,6) {\enspace15\,\%};
\draw[reja, line width=3pt] (axis cs:5.650,7) -- (axis cs:52.420,7);
\draw[tinta, line width=1.2pt] (axis cs:8.197,6.72) -- (axis cs:8.197,7.28);
\fill[serieB, draw=white, line width=0.8pt] (axis cs:52.420,7) circle[radius=3pt];
\fill[serieA, draw=white, line width=0.8pt] (axis cs:5.650,7) circle[radius=3pt];
\node[anchor=west, font=\scriptsize\sffamily\color{tinta2}] at (axis cs:52.420,7) {\enspace52\,\%};
\end{axis}
\end{tikzpicture}

\caption[Falsas alarmas diarias del monitor de Tierra y del adaptado]{Falsas alarmas diarias desde el día 28 en cama.
Naranjo: monitor de Tierra (media del pre-reposo $\pm1{,}96$ desviaciones). Azul: rango adaptado, que por construcción
queda cerca del 5\pct. Raya negra: la misma regla de Tierra aplicada dentro del pre-reposo (placebo), el piso de
falsas alarmas que tiene sin ningún efecto del reposo.}
\label{fig:falsas}
\end{figure}

Parte de esas falsas alarmas no es culpa del reposo. Con 10 a 21 días de pre-reposo, la variación de una persona se
estima mal, y la misma regla aplicada \emph{dentro} del pre-reposo, dejando fuera el día evaluado, ya da entre
\cifra{placebo.min} y \cifra{placebo.max}\pct{} de falsas alarmas según la variable. Ese es el piso. Solo el peso
queda claramente por encima: en el pulso y la presión, el exceso sobre el piso es de pocos puntos (en el pulso,
\cifra{pulso.tierra.fa} contra \cifra{pulso.placebo}\pct) y los intervalos de Tierra incluyen el valor del placebo. En la diuresis y la temperatura, el monitor de Tierra queda incluso
por debajo de su piso.

\begin{table}[htbp]
\caption[Rango esperado: falsas alarmas, ancho y detección por variable]{Rango esperado por variable, desde el día 28
en cama. Falsas alarmas diarias del monitor de Tierra y de su placebo dentro del pre-reposo. Ancho del rango adaptado
dividido por el del monitor de Tierra calibrado, con el mismo 5\pct{} de falsas alarmas. Escalones de 3 días
detectados con la regla de 2 de 3 y diferencia adaptado menos Tierra calibrado, en puntos.}
\label{tab:rango}
\centering\footnotesize\setlength{\tabcolsep}{4.5pt}
\begin{tabular}{lrrrrrr}\toprule
 & \multicolumn{2}{c}{Falsas alarmas diarias (\%)} & \makecell{Ancho} & \multicolumn{3}{c}{Escalones detectados (\%)} \\
\cmidrule(lr){2-3}\cmidrule(lr){5-7}
 Variable & \makecell{Tierra\\{[IC 95\,\%]}} & \makecell{Placebo\\pre-reposo} & \makecell{adaptado /\\Tierra calibr.} & \makecell{Tierra\\calibrado} & Adaptado & \makecell{Diferencia\\{[IC 95\,\%]}} \\
\midrule
Peso & 52 [36;\,67] & 8 & 0,51 & 10 & 26 & +16 [5;\,30] \\
Presión diastólica & 15 [7;\,24] & 9 & 0,77 & 12 & 32 & +19 [9;\,30] \\
Pulso & 14 [7;\,22] & 8 & 0,82 & 11 & 18 & +7 [−1;\,16] \\
Diuresis & 5 [2;\,8] & 10 & 0,85 & 52 & 66 & +15 [0;\,29] \\
Frecuencia respiratoria & 14 [8;\,20] & 12 & 0,89 & 73 & 77 & +4 [0;\,10] \\
Presión sistólica & 11 [7;\,15] & 8 & 0,91 & 44 & 49 & +5 [−4;\,16] \\
Temperatura & 6 [3;\,9] & 9 & 0,94 & 93 & 93 & 0 \\
Agua ingerida\textsuperscript{a} & 5 [2;\,9] & 7 & 0,93 & 92 & 97 & +5 [−2;\,12] \\
\bottomrule\end{tabular}

\end{table}

\begin{decision}
La respuesta a P1 es no, al menos para el peso. Las demás variables tienen un exceso de falsas alarmas sobre el piso
que es chico y que no se distingue bien del ruido. Que el monitor falle sobre todo donde el desplazamiento es grande
(\cref{sec:desplazamiento}) es lo que se espera si el problema es el cambio de nivel.
\end{decision}

\section{P2. Con las mismas falsas alarmas, el adaptado ve más}\label{sec:res-p2}

Comparar el monitor de Tierra con el adaptado tal como están es injusto: el adaptado tiene el ancho calibrado al
5\pct{} y el de Tierra no. La comparación justa es contra el monitor de Tierra \emph{con el ancho calibrado} de la
misma forma, que también queda en torno al 5\pct{} de falsas alarmas diarias (\cifra{prom.calib.fa}\pct{} en
promedio, contra \cifra{prom.adapt.fa}\pct{} del adaptado).

\subsection{Rango más angosto}

Con las mismas falsas alarmas, el rango adaptado es más angosto en todas las variables (\cref{fig:ancho}). En el peso
mide \cifra{peso.ancho.ratio} del de Tierra calibrado: \cifra{peso.adapt.ancho}\un{kg} contra
\cifra{peso.calib.ancho}\un{kg}. En la presión diastólica, \cifra{pa_dia.ancho.ratio}, y en el pulso,
\cifra{pulso.ancho.ratio}. Un rango más angosto con la misma tasa de falsas alarmas quiere decir que el centro está
mejor puesto.

\begin{figure}[htbp]
\centering
\begin{tikzpicture}
\begin{axis}[informe, width=0.78\linewidth, height=6.2cm, xmin=0, xmax=1.1, ymin=-0.6, ymax=7.6, ytick={0,...,7}, yticklabels={Agua,Temperatura,PA sistólica,Respiración,Diuresis,Pulso,PA diastólica,Peso},
 xmajorgrids, xtick={0,0.25,0.5,0.75,1}, xlabel={ancho del rango adaptado / ancho de Tierra calibrado, con las mismas falsas alarmas},
 y axis line style={draw=none}, ytick style={draw=none}, axis x line*=bottom, clip=false]
\draw[tinta3, densely dashed] (axis cs:1,-0.6) -- (axis cs:1,7.6);
\fill[serieA, rounded corners=1.2pt] (axis cs:0,-0.3) rectangle (axis cs:0.9330,0.3);
\node[anchor=west, font=\scriptsize\sffamily\color{tinta2}] at (axis cs:0.9330,0) {0,93};
\fill[serieA, rounded corners=1.2pt] (axis cs:0,0.7) rectangle (axis cs:0.9430,1.3);
\node[anchor=west, font=\scriptsize\sffamily\color{tinta2}] at (axis cs:0.9430,1) {0,94};
\fill[serieA, rounded corners=1.2pt] (axis cs:0,1.7) rectangle (axis cs:0.9070,2.3);
\node[anchor=west, font=\scriptsize\sffamily\color{tinta2}] at (axis cs:0.9070,2) {0,91};
\fill[serieA, rounded corners=1.2pt] (axis cs:0,2.7) rectangle (axis cs:0.8930,3.3);
\node[anchor=west, font=\scriptsize\sffamily\color{tinta2}] at (axis cs:0.8930,3) {0,89};
\fill[serieA, rounded corners=1.2pt] (axis cs:0,3.7) rectangle (axis cs:0.8530,4.3);
\node[anchor=west, font=\scriptsize\sffamily\color{tinta2}] at (axis cs:0.8530,4) {0,85};
\fill[serieA, rounded corners=1.2pt] (axis cs:0,4.7) rectangle (axis cs:0.8240,5.3);
\node[anchor=west, font=\scriptsize\sffamily\color{tinta2}] at (axis cs:0.8240,5) {0,82};
\fill[serieA, rounded corners=1.2pt] (axis cs:0,5.7) rectangle (axis cs:0.7730,6.3);
\node[anchor=west, font=\scriptsize\sffamily\color{tinta2}] at (axis cs:0.7730,6) {0,77};
\fill[serieA, rounded corners=1.2pt] (axis cs:0,6.7) rectangle (axis cs:0.5050,7.3);
\node[anchor=west, font=\scriptsize\sffamily\color{tinta2}] at (axis cs:0.5050,7) {0,51};
\end{axis}
\end{tikzpicture}

\caption[Ancho del rango adaptado frente al de Tierra calibrado]{Ancho del rango adaptado dividido por el del monitor
de Tierra con el ancho calibrado, con el mismo 5\pct{} de falsas alarmas. Bajo 1, el adaptado logra lo mismo con un
rango más angosto.}
\label{fig:ancho}
\end{figure}

\subsection{Más escalones detectados}

Un rango más angosto detecta más. La \cref{fig:deteccion} muestra los escalones sintéticos de 3 días detectados con la
regla de 2 de 3. En el peso, la detección pasa de \cifra{peso.calib.rec.r}\pct{} a \cifra{peso.adapt.rec.r}\pct{}
(\cifra{peso.rec.dif} puntos \cifra{peso.rec.dif.ic}); en la presión diastólica, de \cifra{pa_dia.calib.rec.r} a
\cifra{pa_dia.adapt.rec.r}\pct{} (\cifra{pa_dia.rec.dif} \cifra{pa_dia.rec.dif.ic}), y en el pulso, de
\cifra{pulso.calib.rec.r} a \cifra{pulso.adapt.rec.r}\pct{} (\cifra{pulso.rec.dif} \cifra{pulso.rec.dif.ic}), con
un intervalo que cruza el cero. En la temperatura y el agua los escalones son fáciles de ver con cualquier rango y la
diferencia es mínima.

\begin{figure}[htbp]
\begin{tikzpicture}
\begin{axis}[informe, width=0.86\linewidth, height=7cm, xmin=0, xmax=100, ymin=-0.6, ymax=7.6, ytick={0,...,7}, yticklabels={Agua,Temperatura,PA sistólica,Respiración,Diuresis,Pulso,PA diastólica,Peso},
 xmajorgrids, xtick={0,20,...,100}, xticklabel={\pgfmathprintnumber{\tick}\,\%},
 xlabel={escalones de 3 días detectados, con el ancho de ambos calibrado al 5\,\% de falsas alarmas},
 y axis line style={draw=none}, ytick style={draw=none}, axis x line*=bottom, clip=false,
 legend style={at={(0.02,1.02)}, anchor=south west}, legend columns=2]
\addlegendimage{only marks, mark=*, mark size=2.8pt, serieC}\addlegendentry{Tierra con ancho calibrado\quad}
\addlegendimage{only marks, mark=*, mark size=2.8pt, serieA}\addlegendentry{adaptado}
\draw[reja, line width=3pt] (axis cs:91.800,0) -- (axis cs:96.720,0);
\fill[serieC, draw=white, line width=0.8pt] (axis cs:91.800,0) circle[radius=3pt];
\fill[serieA, draw=white, line width=0.8pt] (axis cs:96.720,0) circle[radius=3pt];
\node[anchor=west, font=\scriptsize\sffamily\color{tinta2}] at (axis cs:96.720,0) {\enspace +5\,pp [−2;\,12]};
\draw[reja, line width=3pt] (axis cs:92.540,1) -- (axis cs:92.540,1);
\fill[serieC, draw=white, line width=0.8pt] (axis cs:92.540,1) circle[radius=3pt];
\fill[serieA, draw=white, line width=0.8pt] (axis cs:92.540,1) circle[radius=3pt];
\node[anchor=west, font=\scriptsize\sffamily\color{tinta2}] at (axis cs:92.540,1) {\enspace sin cambio};
\draw[reja, line width=3pt] (axis cs:43.840,2) -- (axis cs:49.320,2);
\fill[serieC, draw=white, line width=0.8pt] (axis cs:43.840,2) circle[radius=3pt];
\fill[serieA, draw=white, line width=0.8pt] (axis cs:49.320,2) circle[radius=3pt];
\node[anchor=west, font=\scriptsize\sffamily\color{tinta2}] at (axis cs:49.320,2) {\enspace +5\,pp [−4;\,16]};
\draw[reja, line width=3pt] (axis cs:72.600,3) -- (axis cs:76.710,3);
\fill[serieC, draw=white, line width=0.8pt] (axis cs:72.600,3) circle[radius=3pt];
\fill[serieA, draw=white, line width=0.8pt] (axis cs:76.710,3) circle[radius=3pt];
\node[anchor=west, font=\scriptsize\sffamily\color{tinta2}] at (axis cs:76.710,3) {\enspace +4\,pp [0;\,10]};
\draw[reja, line width=3pt] (axis cs:51.610,4) -- (axis cs:66.130,4);
\fill[serieC, draw=white, line width=0.8pt] (axis cs:51.610,4) circle[radius=3pt];
\fill[serieA, draw=white, line width=0.8pt] (axis cs:66.130,4) circle[radius=3pt];
\node[anchor=west, font=\scriptsize\sffamily\color{tinta2}] at (axis cs:66.130,4) {\enspace +15\,pp [0;\,29]};
\draw[reja, line width=3pt] (axis cs:10.960,5) -- (axis cs:17.810,5);
\fill[serieC, draw=white, line width=0.8pt] (axis cs:10.960,5) circle[radius=3pt];
\fill[serieA, draw=white, line width=0.8pt] (axis cs:17.810,5) circle[radius=3pt];
\node[anchor=west, font=\scriptsize\sffamily\color{tinta2}] at (axis cs:17.810,5) {\enspace +7\,pp [−1;\,16]};
\draw[reja, line width=3pt] (axis cs:12.330,6) -- (axis cs:31.510,6);
\fill[serieC, draw=white, line width=0.8pt] (axis cs:12.330,6) circle[radius=3pt];
\fill[serieA, draw=white, line width=0.8pt] (axis cs:31.510,6) circle[radius=3pt];
\node[anchor=west, font=\scriptsize\sffamily\color{tinta2}] at (axis cs:31.510,6) {\enspace +19\,pp [9;\,30]};
\draw[reja, line width=3pt] (axis cs:9.590,7) -- (axis cs:26.030,7);
\fill[serieC, draw=white, line width=0.8pt] (axis cs:9.590,7) circle[radius=3pt];
\fill[serieA, draw=white, line width=0.8pt] (axis cs:26.030,7) circle[radius=3pt];
\node[anchor=west, font=\scriptsize\sffamily\color{tinta2}] at (axis cs:26.030,7) {\enspace +16\,pp [5;\,30]};
\end{axis}
\end{tikzpicture}

\caption[Escalones detectados con las mismas falsas alarmas]{Escalones sintéticos de 3 días detectados desde el día
28, con la regla de 2 de 3 y el ancho de ambos modelos calibrado al 5\pct{} de falsas alarmas. Verde: monitor de
Tierra calibrado. Azul: adaptado con 28 días propios.}
\label{fig:deteccion}
\end{figure}

Promediando sin el agua, el adaptado detecta el \cifra{prom.adapt.rec}\pct{} de los escalones y el monitor de Tierra
calibrado el \cifra{prom.calib.rec}\pct. El monitor de Tierra \emph{sin} calibrar detecta más, el
\cifra{prom.tierra.rec}\pct, pero a cambio de alarmar en el \cifra{prom.tierra.fa}\pct{} de los días normales: con
tantas alarmas, detectar es fácil y no sirve como referencia.

\subsection{A favor y en contra de la deriva}

El signo al azar de los escalones permite ver algo que el promedio esconde. El monitor de Tierra ya está desplazado
respecto del nivel nuevo, así que detecta muy bien lo que va a favor de la deriva y mal lo que va en contra. En el peso,
detecta el \cifra{peso.tierra.recfavor.r}\pct{} de los escalones que lo bajan aún más, pero solo el
\cifra{peso.tierra.reccontra.r}\pct{} de los que lo suben. Calibrado, el desequilibrio se agrava: \cifra{peso.calib.recfavor.r}\pct{}
a favor y \cifra{peso.calib.reccontra.r}\pct{} en contra. El adaptado, centrado en el nivel nuevo, detecta parecido en
los dos sentidos: \cifra{peso.adapt.recfavor.r} y \cifra{peso.adapt.reccontra.r}\pct. Para un monitor de salud eso
importa: un peso que sube de golpe en cama, por ejemplo por retención de líquidos, es justo lo que el monitor de
Tierra no ve.

\subsection{Cuánto es el centro y cuánto el ancho}

El modelo cruzado \modelo{adaptado (ancho de Tierra)} usa el centro adaptado con el ancho de Tierra,
$\pm1{,}96\,\sigma_{s,v}$. Solo con mover el centro, las falsas alarmas diarias del peso bajan de
\cifra{peso.tierra.fa} a \cifra{peso.adapttierra.fa}\pct, y el error del valor esperado baja de
\cifra{peso.tierra.mae} a \cifra{peso.adapt.mae}\un{kg}. Calibrar el ancho hace el resto. En promedio sin el agua, el
centro adaptado lleva las falsas alarmas de \cifra{prom.tierra.fa} a \cifra{prom.adapttierra.fa}\pct{} con el ancho
de Tierra, y el error del valor esperado baja a \cifra{prom.adapt.mae} del de Tierra.

\subsection{Con la regla de 2 de 3}

La \cref{tab:rango} y la \cref{fig:falsas} usan falsas alarmas diarias: basta un día fuera. La \cref{tab:regla} las muestra
también con la regla de 2 de 3, que es la que levanta la alarma. Con la regla, el adaptado da entre
\cifra{adapt.fareg.min} y \cifra{adapt.fareg.max}\pct{} de falsas alarmas según la variable, y el monitor de Tierra
sigue en \cifra{peso.tierra.fareg}\pct{} en el peso.

\begin{table}[htbp]
\caption[Falsas alarmas diarias y con la regla de 2 de 3]{Falsas alarmas (\%) diarias y con la regla de 2 de 3, desde
el día 28 en cama. La última columna es la detección del monitor de Tierra sin calibrar, alta solo porque alarma
mucho.}
\label{tab:regla}
\small
\begin{tabular}{lrrrrrrr}\toprule
 & \multicolumn{3}{c}{Monitor de Tierra} & \makecell{Tierra\\calibrado} & \multicolumn{2}{c}{Adaptado} & \makecell{Detección,\\Tierra} \\
\cmidrule(lr){2-4}\cmidrule(lr){5-5}\cmidrule(lr){6-7}
 Variable & Diaria & 2 de 3 & \makecell{Placebo\\2 de 3} & 2 de 3 & Diaria & 2 de 3 & sin calibrar \\
\midrule
Peso & 52,4 & 53,1 & 6,1 & 5,3 & 5,7 & 4,6 & 55 \\
Presión diastólica & 14,9 & 11,7 & 4,1 & 2,5 & 5,4 & 1,0 & 37 \\
Pulso & 13,7 & 10,9 & 6,3 & 3,1 & 5,1 & 1,2 & 32 \\
Diuresis & 4,7 & 2,5 & 5,2 & 2,8 & 5,1 & 1,7 & 50 \\
Frecuencia respiratoria & 13,8 & 9,6 & 5,1 & 3,9 & 5,0 & 2,5 & 85 \\
Presión sistólica & 11,0 & 6,1 & 3,3 & 1,5 & 4,8 & 0,9 & 62 \\
Temperatura & 6,0 & 3,0 & 4,8 & 1,2 & 5,2 & 1,6 & 90 \\
Agua ingerida & 5,1 & 2,6 & 2,7 & 3,2 & 5,3 & 2,9 & 82 \\
\midrule Promedio sin agua & 16,6 & 13,9 & 5,0 & 2,9 & 5,2 & 1,9 & 58 \\
\bottomrule\end{tabular}

\end{table}

\section{Cuántos días propios hacen falta}\label{sec:curva}

La \cref{fig:curva} separa las dos fuentes de la adaptación. Cada punto es el error del valor esperado dividido por el
error del monitor de Tierra, así que bajo 1 es mejor. La línea verde usa solo los primeros $k$ días propios; la azul
les suma el corrector aprendido de otros sujetos. En $k = 0$, la azul es el corrector solo.

\begin{figure}[htbp]
\centering
\begin{tikzpicture}
\begin{groupplot}[group style={group size=4 by 2, horizontal sep=7mm, vertical sep=15mm, y descriptions at=edge left, x descriptions at=edge bottom}, informe, width=0.295\linewidth, height=4.6cm, xmin=-1.5, xmax=29.5, ymin=0.4, ymax=1.12, xtick={0,7,14,21,28}, ytick={0.4,0.6,0.8,1}, ymajorgrids, xlabel={días propios $k$}, ylabel={error / error de Tierra}, title style={yshift=-1mm}]
\nextgroupplot[title={Peso}, legend to name=leyendacurva, legend columns=2]
\draw[tinta3, densely dashed] (axis cs:-1.5,1) -- (axis cs:29.5,1);
\addplot[serieC, line width=1.6pt, mark=square*, mark size=1.8pt] coordinates {(0,1.0) (3,1.023) (7,0.933) (14,0.784) (21,0.67) (28,0.594)};
\addplot[serieA, line width=1.6pt, mark=*, mark size=2pt] coordinates {(0,0.707) (3,0.63) (7,0.607) (14,0.538) (21,0.493) (28,0.46)};
\addlegendentry{solo días propios\quad}\addlegendentry{días propios + corrector de otros sujetos}
\nextgroupplot[title={PA diastólica}]
\draw[tinta3, densely dashed] (axis cs:-1.5,1) -- (axis cs:29.5,1);
\addplot[serieC, line width=1.6pt, mark=square*, mark size=1.8pt] coordinates {(0,1.0) (3,0.986) (7,0.928) (14,0.858) (21,0.822) (28,0.794)};
\addplot[serieA, line width=1.6pt, mark=*, mark size=2pt] coordinates {(0,1.018) (3,0.943) (7,0.878) (14,0.813) (21,0.789) (28,0.781)};
\nextgroupplot[title={Pulso}]
\draw[tinta3, densely dashed] (axis cs:-1.5,1) -- (axis cs:29.5,1);
\addplot[serieC, line width=1.6pt, mark=square*, mark size=1.8pt] coordinates {(0,1.0) (3,1.08) (7,1.043) (14,1.005) (21,0.932) (28,0.858)};
\addplot[serieA, line width=1.6pt, mark=*, mark size=2pt] coordinates {(0,0.972) (3,0.949) (7,0.901) (14,0.87) (21,0.836) (28,0.8)};
\nextgroupplot[title={Diuresis}]
\draw[tinta3, densely dashed] (axis cs:-1.5,1) -- (axis cs:29.5,1);
\addplot[serieC, line width=1.6pt, mark=square*, mark size=1.8pt] coordinates {(0,1.0) (3,0.852) (7,0.809) (14,0.808) (21,0.81) (28,0.803)};
\addplot[serieA, line width=1.6pt, mark=*, mark size=2pt] coordinates {(0,1.004) (3,0.846) (7,0.823) (14,0.828) (21,0.828) (28,0.821)};
\nextgroupplot[title={Respiración}]
\draw[tinta3, densely dashed] (axis cs:-1.5,1) -- (axis cs:29.5,1);
\addplot[serieC, line width=1.6pt, mark=square*, mark size=1.8pt] coordinates {(0,1.0) (3,0.998) (7,0.978) (14,0.991) (21,0.977) (28,0.963)};
\addplot[serieA, line width=1.6pt, mark=*, mark size=2pt] coordinates {(0,1.014) (3,1.012) (7,0.976) (14,0.989) (21,0.981) (28,0.967)};
\nextgroupplot[title={PA sistólica}]
\draw[tinta3, densely dashed] (axis cs:-1.5,1) -- (axis cs:29.5,1);
\addplot[serieC, line width=1.6pt, mark=square*, mark size=1.8pt] coordinates {(0,1.0) (3,0.965) (7,0.913) (14,0.908) (21,0.887) (28,0.887)};
\addplot[serieA, line width=1.6pt, mark=*, mark size=2pt] coordinates {(0,1.014) (3,0.98) (7,0.924) (14,0.922) (21,0.899) (28,0.898)};
\nextgroupplot[title={Temperatura}]
\draw[tinta3, densely dashed] (axis cs:-1.5,1) -- (axis cs:29.5,1);
\addplot[serieC, line width=1.6pt, mark=square*, mark size=1.8pt] coordinates {(0,1.0) (3,0.966) (7,0.935) (14,0.945) (21,0.921) (28,0.895)};
\addplot[serieA, line width=1.6pt, mark=*, mark size=2pt] coordinates {(0,0.988) (3,0.956) (7,0.919) (14,0.933) (21,0.915) (28,0.887)};
\nextgroupplot[title={Agua}]
\draw[tinta3, densely dashed] (axis cs:-1.5,1) -- (axis cs:29.5,1);
\addplot[serieC, line width=1.6pt, mark=square*, mark size=1.8pt] coordinates {(0,1.0) (3,0.898) (7,0.801) (14,0.802) (21,0.801) (28,0.781)};
\addplot[serieA, line width=1.6pt, mark=*, mark size=2pt] coordinates {(0,0.919) (3,0.898) (7,0.811) (14,0.822) (21,0.824) (28,0.803)};
\end{groupplot}
\node at ($(group c1r2.south west)!0.5!(group c4r2.south east)+(0,-1.45cm)$) {\pgfplotslegendfromname{leyendacurva}};
\end{tikzpicture}

\caption[Curva de aprendizaje del rango esperado]{Error del valor esperado relativo al monitor de Tierra según los días
propios en cama que usa el modelo ($k$). Verde: solo días propios. Azul: días propios más el corrector de otros
sujetos. Se evalúa siempre desde el día 28, así que mide cuánto sirve para predecir lo que viene después, no qué tan
rápido se adapta en las primeras semanas.}
\label{fig:curva}
\end{figure}

En el peso, el corrector solo, sin ningún día propio, ya baja el error a \cifra{curva.peso.pob.0} del de Tierra. Con
28 días propios y sin corrector llega a \cifra{curva.peso.prop.28}, y con los dos a \cifra{curva.peso.pob.28}. En el
pulso y la presión diastólica la curva llega a \cifra{curva.pulso.pob.28} y \cifra{curva.pa_dia.pob.28}.

La \cref{tab:aporte} mide qué agrega el corrector cuando ya están los 28 días propios. Baja el error en más de
\cifra{banda} solo en \cifra{aporte.ayuda}; en el \cifra{aporte.empeora} lo sube, y en las demás variables la
diferencia queda bajo ese umbral. En la mayoría de las variables, entonces, la mejora viene de los días propios. El
umbral de \cifra{banda} es una elección de la auditoría, y estas diferencias no tienen intervalo de confianza.

\begin{table}[htbp]
\caption[Qué aporta el corrector de otros sujetos]{Error del valor esperado relativo al monitor de Tierra. El aporte del
corrector es la diferencia entre usar ambos y usar solo los días propios, con $k = 28$; se clasifica con una banda de
$\pm$\cifra{banda}.}
\label{tab:aporte}
\small
\begin{tabular}{lrrrrl}\toprule
 & \makecell{Solo corrector\\($k=0$)} & \makecell{Solo días propios\\($k=28$)} & \makecell{Ambos\\($k=28$)} & \makecell{Ambos $-$\\propios} & \makecell[l]{Aporte del\\corrector} \\
\midrule
Peso & 0,71 & 0,59 & 0,46 & −0,13 & mejora \\
Presión diastólica & 1,02 & 0,79 & 0,78 & −0,01 & sin diferencia \\
Pulso & 0,97 & 0,86 & 0,80 & −0,06 & mejora \\
Diuresis & 1,00 & 0,80 & 0,82 & +0,02 & sin diferencia \\
Frecuencia respiratoria & 1,01 & 0,96 & 0,97 & +0,00 & sin diferencia \\
Presión sistólica & 1,01 & 0,89 & 0,90 & +0,01 & sin diferencia \\
Temperatura & 0,99 & 0,90 & 0,89 & −0,01 & sin diferencia \\
Agua ingerida & 0,92 & 0,78 & 0,80 & +0,02 & empeora \\
\bottomrule\end{tabular}

\end{table}

\section{Persona por persona}\label{sec:personas}

El ancho calibrado es uno para toda la población, así que el 5\pct{} de falsas alarmas se cumple en promedio, no en
cada persona. La \cref{fig:fa-persona} lo muestra en el peso: de \cifra{fapeso.n} personas, \cifra{fapeso.cero} no
tienen ninguna falsa alarma y \cifra{fapeso.sobre10} pasan del 10\pct, con un máximo de \cifra{fapeso.max}\pct. La
mediana es \cifra{fapeso.mediana}\pct{} y la media \cifra{fapeso.media}\pct, con entre \cifra{fapeso.dias_min} y
\cifra{fapeso.dias_max} días evaluados por persona.

\begin{figure}[htbp]
\begin{tikzpicture}
\begin{axis}[informe, width=0.92\linewidth, height=5.6cm, xmin=0, xmax=30, ymin=0, ymax=45, ymajorgrids, ytick={0,10,...,40}, yticklabel={\pgfmathprintnumber{\tick}\,\%}, xtick=\empty, xlabel={cada barra es una persona, ordenadas de más a menos falsas alarmas}, ylabel={falsas alarmas diarias}, axis x line*=bottom, axis y line*=left, clip=false]
\draw[serieB, densely dashed, line width=1pt] (axis cs:0,6.2) -- (axis cs:30,6.2) node[pos=1, anchor=south east, font=\scriptsize\sffamily\color{tinta2}] {promedio de la población: 6,2\,\%};
\fill[serieA, rounded corners=0.8pt] (axis cs:0.6599999999999999,0) rectangle (axis cs:1.34,40.7);
\fill[serieA, rounded corners=0.8pt] (axis cs:1.66,0) rectangle (axis cs:2.34,38.1);
\fill[serieA, rounded corners=0.8pt] (axis cs:2.66,0) rectangle (axis cs:3.34,34.3);
\fill[serieA, rounded corners=0.8pt] (axis cs:3.66,0) rectangle (axis cs:4.34,21.4);
\fill[serieA, rounded corners=0.8pt] (axis cs:4.66,0) rectangle (axis cs:5.34,18.8);
\fill[serieA, rounded corners=0.8pt] (axis cs:5.66,0) rectangle (axis cs:6.34,14.3);
\fill[serieA, rounded corners=0.8pt] (axis cs:6.66,0) rectangle (axis cs:7.34,3.2);
\fill[serieA, rounded corners=0.8pt] (axis cs:7.66,0) rectangle (axis cs:8.34,3.2);
\fill[serieA, rounded corners=0.8pt] (axis cs:8.66,0) rectangle (axis cs:9.34,3.1);
\fill[serieA, rounded corners=0.8pt] (axis cs:9.66,0) rectangle (axis cs:10.34,1.6);
\fill[serieA, rounded corners=0.8pt] (axis cs:10.66,0) rectangle (axis cs:11.34,1.6);
\fill[tinta3] (axis cs:12,0.0) circle[radius=1.4pt];
\fill[tinta3] (axis cs:13,0.0) circle[radius=1.4pt];
\fill[tinta3] (axis cs:14,0.0) circle[radius=1.4pt];
\fill[tinta3] (axis cs:15,0.0) circle[radius=1.4pt];
\fill[tinta3] (axis cs:16,0.0) circle[radius=1.4pt];
\fill[tinta3] (axis cs:17,0.0) circle[radius=1.4pt];
\fill[tinta3] (axis cs:18,0.0) circle[radius=1.4pt];
\fill[tinta3] (axis cs:19,0.0) circle[radius=1.4pt];
\fill[tinta3] (axis cs:20,0.0) circle[radius=1.4pt];
\fill[tinta3] (axis cs:21,0.0) circle[radius=1.4pt];
\fill[tinta3] (axis cs:22,0.0) circle[radius=1.4pt];
\fill[tinta3] (axis cs:23,0.0) circle[radius=1.4pt];
\fill[tinta3] (axis cs:24,0.0) circle[radius=1.4pt];
\fill[tinta3] (axis cs:25,0.0) circle[radius=1.4pt];
\fill[tinta3] (axis cs:26,0.0) circle[radius=1.4pt];
\fill[tinta3] (axis cs:27,0.0) circle[radius=1.4pt];
\fill[tinta3] (axis cs:28,0.0) circle[radius=1.4pt];
\fill[tinta3] (axis cs:29,0.0) circle[radius=1.4pt];
\node[anchor=south, font=\scriptsize\sffamily\color{tinta2}, align=center] at (axis cs:20.5,13) {18 de 29 personas\\sin ninguna falsa alarma};
\end{axis}
\end{tikzpicture}

\caption[Falsas alarmas del adaptado persona por persona]{Falsas alarmas diarias del rango adaptado en el peso,
persona por persona, ordenadas de más a menos. La línea punteada es el promedio de la población.}
\label{fig:fa-persona}
\end{figure}

El error del valor esperado también se reparte de forma desigual (\cref{tab:por-persona}). En el peso, el adaptado se
equivoca menos que el monitor de Tierra en \cifra{porpersona.peso.mejoran} de las \cifra{porpersona.peso.n} personas, y la mediana del
cociente de error es \cifra{porpersona.peso.mediana}; pero en la peor persona el cociente llega a \cifra{porpersona.peso.peor}.

\begin{table}[htbp]
\caption[Error del adaptado frente al de Tierra, por persona]{Error del valor esperado del adaptado ($k = 28$) dividido
por el del monitor de Tierra, calculado persona por persona.}
\label{tab:por-persona}
\small
\begin{tabular}{lrrrr}\toprule
 Variable & Personas & \makecell{Mejoran con\\el adaptado} & \makecell{Mediana del\\cociente de error} & \makecell{Peor\\cociente} \\
\midrule
Peso & 29 & 20 & 0,36 & 3,58 \\
Presión diastólica & 29 & 21 & 0,92 & 1,28 \\
Pulso & 29 & 23 & 0,86 & 1,44 \\
Diuresis & 24 & 19 & 0,85 & 1,43 \\
Frecuencia respiratoria & 29 & 16 & 1,00 & 2,55 \\
Presión sistólica & 29 & 19 & 0,95 & 1,27 \\
Temperatura & 26 & 16 & 0,96 & 1,18 \\
Agua ingerida & 23 & 15 & 0,96 & 1,62 \\
\bottomrule\end{tabular}

\end{table}

\section{P3. El pronóstico a 7 días}\label{sec:res-pronostico}

La \cref{tab:pronostico} y la \cref{fig:pronostico} muestran el MAE relativo a la persistencia, promediado sin el agua. En
cama, todos los PatchTST le ganan a la media de los últimos 7 días (\cifra{pron.media7}), y Chronos-2 sin adaptar
queda detrás (\cifra{pron.chronos}). El mejor es el ajuste completo desde Tierra, con \cifra{pron.todo}, pero los
cinco PatchTST quedan muy juntos, entre \cifra{pron.todo} y \cifra{pron.tierra}.

\begin{figure}[htbp]
\centering
\begin{tikzpicture}
\begin{axis}[informe, width=0.62\linewidth, height=6cm, xmin=0.785, xmax=0.825, ymin=-0.6, ymax=6.6, ytick={0,...,6}, yticklabels={{Chronos-2 sin adaptar},{Media de los últimos 7 días},{Tierra sin adaptar},{Solo reposo, desde cero},{Pre-reposo + reposo, desde cero (sin LifeSnaps)},{Tierra → espacio, solo la cabeza},{Tierra → espacio, ajuste completo}},
 xmajorgrids, xtick={0.79,0.80,0.81,0.82}, x tick label style={/pgf/number format/fixed, /pgf/number format/precision=2},
 xlabel={MAE relativo a persistencia, promedio sin agua (menos es mejor)},
 y axis line style={draw=none}, ytick style={draw=none}, axis x line*=bottom, clip=false]
\draw[reja, line width=1pt] (axis cs:0.785,0) -- (axis cs:0.8200,0);
\fill[tinta3, draw=white, line width=0.8pt] (axis cs:0.8200,0) circle[radius=3.2pt];
\node[anchor=west, font=\scriptsize\sffamily\color{tinta2}] at (axis cs:0.8200,0) {\enspace0,820};
\draw[reja, line width=1pt] (axis cs:0.785,1) -- (axis cs:0.8070,1);
\fill[tinta3, draw=white, line width=0.8pt] (axis cs:0.8070,1) circle[radius=3.2pt];
\node[anchor=west, font=\scriptsize\sffamily\color{tinta2}] at (axis cs:0.8070,1) {\enspace0,807};
\draw[reja, line width=1pt] (axis cs:0.785,2) -- (axis cs:0.7990,2);
\fill[serieB, draw=white, line width=0.8pt] (axis cs:0.7990,2) circle[radius=3.2pt];
\node[anchor=west, font=\scriptsize\sffamily\color{tinta2}] at (axis cs:0.7990,2) {\enspace0,799};
\draw[reja, line width=1pt] (axis cs:0.785,3) -- (axis cs:0.7970,3);
\fill[tinta2, draw=white, line width=0.8pt] (axis cs:0.7970,3) circle[radius=3.2pt];
\node[anchor=west, font=\scriptsize\sffamily\color{tinta2}] at (axis cs:0.7970,3) {\enspace0,797};
\draw[reja, line width=1pt] (axis cs:0.785,4) -- (axis cs:0.7960,4);
\fill[tinta2, draw=white, line width=0.8pt] (axis cs:0.7960,4) circle[radius=3.2pt];
\node[anchor=west, font=\scriptsize\sffamily\color{tinta2}] at (axis cs:0.7960,4) {\enspace0,796};
\draw[reja, line width=1pt] (axis cs:0.785,5) -- (axis cs:0.7960,5);
\fill[serieA, draw=white, line width=0.8pt] (axis cs:0.7960,5) circle[radius=3.2pt];
\node[anchor=west, font=\scriptsize\sffamily\color{tinta2}] at (axis cs:0.7960,5) {\enspace0,796};
\draw[reja, line width=1pt] (axis cs:0.785,6) -- (axis cs:0.7930,6);
\fill[serieA, draw=white, line width=0.8pt] (axis cs:0.7930,6) circle[radius=3.2pt];
\node[anchor=west, font=\scriptsize\sffamily\color{tinta2}] at (axis cs:0.7930,6) {\enspace0,793};
\end{axis}
\end{tikzpicture}

\caption[MAE relativo a persistencia de cada modelo]{Pronóstico a 7 días en cama: MAE relativo a la persistencia,
promedio de las siete variables sin el agua. Naranjo: modelo de Tierra sin adaptar. Azul: Tierra adaptado al reposo.
Gris: referencias y modelos desde cero.}
\label{fig:pronostico}
\end{figure}

\begin{table}[htbp]
\caption[Pronóstico a 7 días por modelo]{MAE relativo a la persistencia, promedio sin el agua, en cama y en el
pre-reposo de los sujetos de prueba. Los modelos adaptados al reposo y los entrenados desde cero en el reposo solo se
evalúan en cama.}
\label{tab:pronostico}
\small
\begin{tabular}{lrr}\toprule
 Modelo & \makecell{En cama} & \makecell{Pre-reposo} \\
\midrule
Chronos-2 sin adaptar & 0,820 & 0,993 \\
Media de los últimos 7 días & 0,807 & 0,898 \\
Tierra sin adaptar & 0,799 & 0,891 \\
Solo reposo, desde cero & 0,797 & — \\
Pre-reposo + reposo, desde cero (sin LifeSnaps) & 0,796 & — \\
Tierra → espacio, solo la cabeza & 0,796 & — \\
Tierra → espacio, ajuste completo & 0,793 & — \\
\bottomrule\end{tabular}

\end{table}

Las diferencias son de milésimas, así que lo que importa son los pares, con su intervalo (\cref{fig:pares}):
\begin{description}
\item[Contra la media de 7 días] Tierra → espacio le gana por \cifra{par.cabeza.media7.r} \cifra{par.cabeza.media7.icr}
  ajustando solo la cabeza y por \cifra{par.todo.media7.r} \cifra{par.todo.media7.icr} ajustando todo. La mejora es
  real pero chica: \cifra{par.cabeza.media7.puntos} y \cifra{par.todo.media7.puntos} puntos del error de la
  persistencia.
\item[Contra los modelos desde cero] Ajustando solo la cabeza es un empate: \cifra{par.cabeza.cero}
  \cifra{par.cabeza.cero.ic} contra el reposo desde cero y \cifra{par.cabeza.sinls} \cifra{par.cabeza.sinls.ic}
  contra el pre-reposo + reposo sin LifeSnaps. Con ajuste completo gana por poco: \cifra{par.todo.cero}
  \cifra{par.todo.cero.ic} y \cifra{par.todo.sinls} \cifra{par.todo.sinls.ic}.
\item[Contra Tierra sin adaptar] El ajuste completo mejora al modelo de Tierra en \cifra{par.todo.tierra}
  \cifra{par.todo.tierra.ic}; solo la cabeza, en \cifra{par.cabeza.tierra} \cifra{par.cabeza.tierra.ic}, con un
  intervalo que toca el cero.
\item[Chronos-2] Sin adaptar le va peor que a la media de 7 días: \cifra{par.chronos.media7.r}
  \cifra{par.chronos.media7.icr}. Recibe el contexto con los huecos rellenados, así que es solo una referencia.
\end{description}

\begin{figure}[htbp]
\centering
\begin{tikzpicture}
\begin{axis}[informe, width=0.42\linewidth, height=8.4cm, xmin=-0.022, xmax=0.022, ymin=-0.6, ymax=9.6, ytick={0,...,9}, yticklabels={{ajuste completo − solo la cabeza},{solo la cabeza − Tierra sin adaptar},{ajuste completo − Tierra sin adaptar},{ajuste completo − pre-reposo + reposo sin LifeSnaps},{ajuste completo − reposo desde cero},{Chronos-2 − media de 7 días},{solo la cabeza − pre-reposo + reposo sin LifeSnaps},{solo la cabeza − reposo desde cero},{ajuste completo − media de 7 días},{solo la cabeza − media de 7 días}},
 xmajorgrids, xtick={-0.02,-0.01,0,0.01,0.02}, scaled x ticks=false,
 x tick label style={/pgf/number format/fixed, /pgf/number format/precision=2},
 xlabel={diferencia de MAE relativo (negativo: gana el primero)},
 y axis line style={draw=none}, ytick style={draw=none}, axis x line*=bottom, clip=false,
 legend style={at={(0.5,1.02)}, anchor=south}, legend columns=2]
\addlegendimage{only marks, mark=*, mark size=2.6pt, tinta}\addlegendentry{informado por el kernel\quad}
\addlegendimage{only marks, mark=*, mark size=2.6pt, mark options={fill=white, draw=serieA, line width=1pt}}\addlegendentry{recalculado en la auditoría}
\draw[tinta3, densely dashed] (axis cs:0,-0.6) -- (axis cs:0,9.6);
\draw[serieA, line width=1.4pt] (axis cs:-0.0049,0) -- (axis cs:-0.0011,0);
\filldraw[fill=white, draw=serieA, line width=1.1pt] (axis cs:-0.003,0) circle[radius=2.6pt];
\node[anchor=west, font=\scriptsize\sffamily\color{tinta2}] at (axis cs:0.0225,0) {−0,0030 [−0,0049;\,−0,0011]};
\draw[serieA, line width=1.4pt] (axis cs:-0.0066,1) -- (axis cs:0.0001,1);
\filldraw[fill=white, draw=serieA, line width=1.1pt] (axis cs:-0.0033,1) circle[radius=2.6pt];
\node[anchor=west, font=\scriptsize\sffamily\color{tinta2}] at (axis cs:0.0225,1) {−0,0033 [−0,0066;\,0,0001]};
\draw[serieA, line width=1.4pt] (axis cs:-0.009,2) -- (axis cs:-0.0037,2);
\filldraw[fill=white, draw=serieA, line width=1.1pt] (axis cs:-0.0062,2) circle[radius=2.6pt];
\node[anchor=west, font=\scriptsize\sffamily\color{tinta2}] at (axis cs:0.0225,2) {−0,0062 [−0,0090;\,−0,0037]};
\draw[serieA, line width=1.4pt] (axis cs:-0.0047,3) -- (axis cs:-0.0007,3);
\filldraw[fill=white, draw=serieA, line width=1.1pt] (axis cs:-0.0027,3) circle[radius=2.6pt];
\node[anchor=west, font=\scriptsize\sffamily\color{tinta2}] at (axis cs:0.0225,3) {−0,0027 [−0,0047;\,−0,0007]};
\draw[serieA, line width=1.4pt] (axis cs:-0.0064,4) -- (axis cs:-0.0011,4);
\filldraw[fill=white, draw=serieA, line width=1.1pt] (axis cs:-0.0038,4) circle[radius=2.6pt];
\node[anchor=west, font=\scriptsize\sffamily\color{tinta2}] at (axis cs:0.0225,4) {−0,0038 [−0,0064;\,−0,0011]};
\draw[tinta, line width=1.4pt] (axis cs:0.007,5) -- (axis cs:0.019,5);
\fill[tinta, draw=white, line width=0.8pt] (axis cs:0.0129,5) circle[radius=2.8pt];
\node[anchor=west, font=\scriptsize\sffamily\color{tinta2}] at (axis cs:0.0225,5) {+0,0129 [0,0070;\,0,0190]};
\draw[tinta, line width=1.4pt] (axis cs:-0.002,6) -- (axis cs:0.0027,6);
\fill[tinta, draw=white, line width=0.8pt] (axis cs:0.0003,6) circle[radius=2.8pt];
\node[anchor=west, font=\scriptsize\sffamily\color{tinta2}] at (axis cs:0.0225,6) {+0,0003 [−0,0020;\,0,0027]};
\draw[tinta, line width=1.4pt] (axis cs:-0.0038,7) -- (axis cs:0.002,7);
\fill[tinta, draw=white, line width=0.8pt] (axis cs:-0.0008,7) circle[radius=2.8pt];
\node[anchor=west, font=\scriptsize\sffamily\color{tinta2}] at (axis cs:0.0225,7) {−0,0008 [−0,0038;\,0,0020]};
\draw[tinta, line width=1.4pt] (axis cs:-0.0193,8) -- (axis cs:-0.0077,8);
\fill[tinta, draw=white, line width=0.8pt] (axis cs:-0.0137,8) circle[radius=2.8pt];
\node[anchor=west, font=\scriptsize\sffamily\color{tinta2}] at (axis cs:0.0225,8) {−0,0137 [−0,0193;\,−0,0077]};
\draw[tinta, line width=1.4pt] (axis cs:-0.0163,9) -- (axis cs:-0.0045,9);
\fill[tinta, draw=white, line width=0.8pt] (axis cs:-0.0107,9) circle[radius=2.8pt];
\node[anchor=west, font=\scriptsize\sffamily\color{tinta2}] at (axis cs:0.0225,9) {−0,0107 [−0,0163;\,−0,0045]};
\end{axis}
\end{tikzpicture}

\caption[Diferencias pareadas del pronóstico]{Diferencias pareadas por sujeto del MAE relativo a la persistencia, con
IC~95\pct{} por bootstrap de sujetos. Bajo 0, el primer modelo se equivoca menos. Punto lleno: par que informó el
kernel; punto hueco: par recalculado en la auditoría con la misma función del proyecto.}
\label{fig:pares}
\end{figure}

En el pre-reposo, donde no hay que adaptar nada, el modelo de Tierra sin adaptar queda en \cifra{pron.pre.tierra},
apenas mejor que la media de 7 días (\cifra{pron.pre.media7}), y Chronos-2 en \cifra{pron.pre.chronos}, casi igual que
la persistencia. El detalle por variable está en la \cref{tab:pronostico-var}.

\begin{cuidado}
Con una sola semilla y brazos entrenados en condiciones distintas (\cref{sec:entreno}), una ventaja de
\cifra{par.todo.cero.r} o menos no alcanza para decir si el pre-entrenamiento en Tierra aporta. Tampoco alcanza para decir que no
aporta. La corrida v4 iguala las condiciones y usa tres semillas (\cref{sec:v4}).
\end{cuidado}

\section{Cómo se entrenaron las redes}\label{sec:entreno}

La \cref{tab:entreno} muestra en qué época cada brazo alcanzó su mejor pérdida de validación. Los modelos desde cero
llegan a su mejor época entre la \cifra{cero.ep.min} y la \cifra{cero.ep.max}, de un tope de 60: con una tasa de
aprendizaje de \texttt{1e-3} sobreajustan pronto. El ajuste completo, con una tasa cinco veces menor y partiendo de un
modelo ya entrenado, también para temprano. \enquote{Solo la cabeza} es el caso contrario: el tope de 30 épocas lo
cortó en \cifra{cabeza.cortados} de 5 folds, y en el fold~0 su mejor época fue la última.

\begin{table}[htbp]
\caption[Épocas de entrenamiento por brazo y fold]{Mejor época de cada brazo y fold (contada desde 1). En negrita, una
mejor época igual al tope. \enquote{Cortados por el tope}: folds en que el tope llegó antes de que la paciencia de 6
épocas pudiera detener el entrenamiento.}
\label{tab:entreno}
\small
\begin{tabular}{lcrrrrrrr}\toprule
 & & & \multicolumn{5}{c}{Mejor época por fold} & \\
\cmidrule(lr){4-8}
 Brazo & lr & Tope & 0 & 1 & 2 & 3 & 4 & \makecell{Cortados\\por el tope} \\
\midrule
Base de Tierra & \texttt{1e-3} & 60 & 12 & 17 & 17 & 15 & 11 & 0 \\
Reposo, desde cero & \texttt{1e-3} & 60 & 5 & 3 & 5 & 5 & 5 & 0 \\
Pre-reposo + reposo, desde cero & \texttt{1e-3} & 60 & 17 & 4 & 7 & 10 & 6 & 0 \\
Tierra → espacio, solo la cabeza & \texttt{1e-3} & 30 & \textbf{30} & 28 & 25 & 15 & 5 & 3 \\
Tierra → espacio, ajuste completo & \texttt{2e-4} & 30 & 15 & 7 & 4 & 3 & 2 & 0 \\
\bottomrule\end{tabular}

\end{table}

La \cref{fig:historia} muestra la pérdida de validación del fold~0. La de \enquote{solo la cabeza} baja de forma lenta y
sostenida, de \cifra{hist.f0.cabeza.prim} a \cifra{hist.f0.cabeza.ult}, y seguía bajando cuando se cortó. El modelo
desde cero sobreajusta en pocas épocas.

\begin{figure}[htbp]
\centering
\begin{tikzpicture}
\begin{axis}[informe, width=0.82\linewidth, height=6.2cm, xmin=0, xmax=33, ymajorgrids, xlabel={época}, ylabel={pérdida de validación (MSE escalado)}, axis x line*=bottom, axis y line*=left, clip=false, xtick={0,5,...,30}, y tick label style={/pgf/number format/fixed, /pgf/number format/precision=2}]
\addplot[tinta3, line width=1.5pt, mark=*, mark size=1.2pt] coordinates {(1,3.17139) (2,3.16791) (3,3.12855) (4,3.08697) (5,3.08573) (6,3.11419) (7,3.14897) (8,3.25872) (9,3.16860) (10,3.09800) (11,3.54830)};
\node[anchor=west, font=\scriptsize\sffamily\color{tinta2}] at (axis cs:11,3.54830) {\, reposo desde cero};
\addplot[serieC, line width=1.5pt, mark=*, mark size=1.2pt] coordinates {(1,3.16873) (2,3.16085) (3,3.15889) (4,3.15214) (5,3.14798) (6,3.14439) (7,3.14352) (8,3.13900) (9,3.13466) (10,3.13637) (11,3.13316) (12,3.13074) (13,3.12544) (14,3.12813) (15,3.13083) (16,3.12448) (17,3.12500) (18,3.12240) (19,3.12289) (20,3.12606) (21,3.12339) (22,3.12191) (23,3.12276) (24,3.12456) (25,3.12180) (26,3.12556) (27,3.12235) (28,3.12046) (29,3.12190) (30,3.11946)};
\node[anchor=west, font=\scriptsize\sffamily\color{tinta2}] at (axis cs:30,3.11946) {\, solo la cabeza};
\addplot[serieA, line width=1.5pt, mark=*, mark size=1.2pt] coordinates {(1,3.10849) (2,3.10577) (3,3.10536) (4,3.09487) (5,3.09256) (6,3.08953) (7,3.07840) (8,3.07787) (9,3.08025) (10,3.07249) (11,3.06811) (12,3.06254) (13,3.06953) (14,3.06265) (15,3.05707) (16,3.05822) (17,3.07858) (18,3.07381) (19,3.07577) (20,3.07000) (21,3.08060)};
\node[anchor=west, font=\scriptsize\sffamily\color{tinta2}] at (axis cs:21,3.08060) {\, ajuste completo};
\draw[tinta3, densely dashed] (axis cs:30,\pgfkeysvalueof{/pgfplots/ymin}) -- (axis cs:30,\pgfkeysvalueof{/pgfplots/ymax});
\end{axis}
\end{tikzpicture}

\caption[Pérdida de validación en el fold 0]{Pérdida de validación del fold~0 por época. La línea punteada marca el
tope de 30 épocas de los brazos de adaptación.}
\label{fig:historia}
\end{figure}

\section{Qué cambia en los pesos}\label{sec:pesos}

El modelo publicado \texttt{nova-espacio} es \texttt{nova-tierra} adaptado en modo \enquote{solo la cabeza}. Si ese
modo congelara todo el backbone, entre los dos solo cambiarían los tensores de la cabeza. Comparando los archivos
\texttt{.safetensors}, cambian \cifra{pesos.cambian} de \cifra{pesos.tensores}: los \cifra{pesos.cabeza} de la cabeza
y \cifra{pesos.bn} de BatchNorm (\cref{tab:pesos}). Ningún peso entrenable del backbone cambia; cambian las medias y
varianzas móviles de BatchNorm y su contador de lotes, que pasa de \cifra{pesos.nb.tierra} a \cifra{pesos.nb.espacio}.

\begin{table}[htbp]
\caption[Tensores que cambian entre los pesos publicados]{Tensores que cambian entre \texttt{nova-tierra} y
\texttt{nova-espacio}.}
\label{tab:pesos}
\small
\begin{tabularx}{\textwidth}{@{}>{\raggedright\arraybackslash}X l r@{}}
\toprule
Tensores & Tipo & Cuántos \\
\midrule
\texttt{head.projection.weight}, \texttt{head.projection.bias} & parámetros de la cabeza & \cifra{pesos.cabeza} \\
\texttt{encoder.layers.\{0,1,2\}.norm\_sublayer\{1,3\}.batchnorm}: \texttt{running\_mean}, \texttt{running\_var},
\texttt{num\_batches\_tracked} & buffers de BatchNorm & \cifra{pesos.bn} \\
\midrule
Total & & \cifra{pesos.cambian} de \cifra{pesos.tensores} \\
\bottomrule
\end{tabularx}
\end{table}

La causa está en el código: congelar el backbone solo apaga sus gradientes, pero el entrenamiento pone el modelo en
modo de entrenamiento, y en ese modo BatchNorm actualiza sus estadísticas con cada lote (ecuación~\ref{eq:bn}) y el dropout
sigue activo. El brazo \enquote{solo la cabeza} es, en la práctica, la cabeza más una normalización re-estimada con los
datos del reposo, una forma de AdaBN \parencite{li2018adabn}. Los folds usan la misma función, así que lo mismo pasa
en las cifras de la \cref{sec:res-pronostico}. Ninguna cifra de desempeño sale de los pesos publicados; de ellos sale solo
este conteo.
